\documentclass[12pt]{article}
% Shared print layout. The web's provisional design and MathJax are independent.
\usepackage[utf8]{inputenc}
\usepackage[T1]{fontenc}
\usepackage[spanish,es-nodecimaldot]{babel}
\usepackage{lmodern}
\usepackage{amsmath,amssymb}
\usepackage{graphicx}
\usepackage{geometry}
\geometry{a4paper,left=20mm,right=20mm,top=25mm,bottom=22mm,headheight=16pt,headsep=8mm}
\usepackage{fancyhdr,lastpage}
\usepackage{enumitem}
\usepackage{xcolor}
\usepackage{pgfplots}
\pgfplotsset{compat=1.18}
\usepgfplotslibrary{fillbetween}
\usepackage[hidelinks,pdfencoding=auto]{hyperref}
\hypersetup{pdflang=es-ES,pdfauthor={},pdfsubject={Matemáticas Aplicadas a las Ciencias Sociales II}}
\pagestyle{fancy}
\fancyhf{}
\fancyhead[L]{\small Matemáticas Aplicadas a las Ciencias Sociales II}
\fancyhead[R]{\small 2.\textsuperscript{o} Bachillerato}
\fancyfoot[L]{\footnotesize Subida de nota T1 | Edición revisada 1.0}
\fancyfoot[R]{\footnotesize Página \thepage\ de \pageref{LastPage}}
\renewcommand{\headrulewidth}{0.4pt}
\setlength{\parindent}{0pt}
\setlength{\parskip}{0.5em}
\setlist{itemsep=0.3em,topsep=0.3em}
\newcommand{\apartado}[1]{\subsection*{#1}}
\pgfplotsset{
  revision/.style={
    width=0.94\linewidth,height=7.5cm,
    axis lines=middle,grid=major,grid style={gray!18},
    tick label style={font=\small},label style={font=\small},
    xlabel={$x$},ylabel={$y$},samples=160,
    legend style={at={(0.5,-0.16)},anchor=north,draw=none,font=\small},
    clip=true
  }
}

\hypersetup{pdftitle={Subida de nota T1 - Examen revisado}}
\begin{document}
\section*{Subida de nota T1: examen general}
\textbf{Matemáticas Aplicadas a las Ciencias Sociales II}\\
Duración: \textbf{55 minutos}. Puntuación total: \textbf{10 puntos}.

Apellidos y nombre: \hrulefill

Fecha: \hrulefill

Lee todas las preguntas y responde de forma razonada. Numera las páginas de tu respuesta. Se permite el uso de calculadora.
\section*{Ejercicio 1 \hfill 4 puntos}
\textit{Tiempo orientativo: 15 minutos.} Se considera la función
\[
f(x)=\begin{cases}\dfrac{a}{x-1},&x<0,\\x^2-bx-1,&x\geq0.\end{cases}
\]
\begin{enumerate}[label=(\alph*)]
\item (2 puntos) Calcula $a$ y $b$ para que $f$ sea derivable en $x=0$.
\item (2 puntos) Para $a=1$ y $b=2$, calcula el área comprendida entre la gráfica de $f$, el eje de abscisas y las rectas $x=0$ y $x=1$.
\end{enumerate}
\section*{Ejercicio 2 \hfill 5 puntos}
\textit{Tiempo orientativo: 35 minutos.} Dada la función
\[g(x)=\frac{x}{x+3},\]
\begin{enumerate}[label=(\alph*)]
\item (3,5 puntos) Estudia continuidad, asíntotas, monotonía y puntos de corte de $g$, y represéntala gráficamente.
\item (1,5 puntos) Halla la ecuación de la recta tangente a la gráfica de $g$ en el punto de abscisa $x=2$.
\end{enumerate}
\section*{Ejercicio 3 \hfill 1 punto}
\textit{Tiempo orientativo: 5 minutos.} Calcula una primitiva de
\[h(x)=x^2+e^{2x}-\frac{1}{1-x}.\]
\end{document}
